\exercisesection{DIMENSION OF NOETHERIAN RINGS}

\begin{enumerate}
\def\labelenumi{\arabic{enumi})}
\tightlist
\item
  Let A be a ring, and let U be an open subset of Spec(A).
\end{enumerate}

\begin{enumerate}
\def\labelenumi{\alph{enumi})}
\item
  Show that the map \(\mathfrak{p} \mapsto \mathrm{V}(\mathfrak{p}) \cap \mathrm{U}\) induces a bijection from U onto the set of closed irreducible subsets of U.
\item
  Show that dim(U) is the upper bound of the set of lengths of chains of prime ideals of A belonging to U.
\item
  Let \(\Phi\) be the set of minimal prime ideals of A belonging to U. Show that one has
\end{enumerate}

\[
\dim (\mathrm{U}) = \sup _ {\mathfrak {p} \in \Phi} \dim (\mathrm{V} (\mathfrak {p}) \cap \mathrm{U}).
\]

\begin{enumerate}
\def\labelenumi{\alph{enumi})}
\setcounter{enumi}{3}
\item
  Assume henceforth that A is local and U ≠ Spec(A). If A is Noetherian, establish the equality \(\dim(\mathrm{U}) = \sup_{\mathfrak{p}\in\Phi}(\dim(\mathrm{A}/\mathfrak{p}) - 1)\).
\item
  Let h be an integer such that 0 \textless{} h \textless{} dim(A). Let \(E_{h}\) be the set of prime ideals \(p \in U\) such that \(ht(p) = h\) and \(\dim(A/p) = \dim(A) - h\) . If A is Noetherian, show that \(E_{h}\) is infinite. Show that \(E_{h}\) may be finite when A is not Noetherian (take for A a valuation ring of height 2, for U the open \(\text{Spec}(A) - \{m_{A}\}\) , and choose h = 1).
\end{enumerate}

\begin{enumerate}
\def\labelenumi{\arabic{enumi})}
\setcounter{enumi}{1}
\item
  Let A be a Noetherian integral domain. Suppose that A is not a field and that the field of fractions of A is a finitely generated A-algebra. Then A is semi-local and of dimension 1. (Show that there exists a nonzero, non-invertible element \(f\) of A such that \(A[f^{-1}]\) is a field, and reduce to proving that \(\dim(A/(f)) = 0\) ; if it were otherwise, \(A/(f)\) would possess a prime ideal of height 1 whose inverse image p in A would be a prime ideal of height 2; but the local ring \(A_p\) of dimension 2 could have only a finite number of prime ideals, contrary to prop. 6 of § 3, no. 3.)
\item
  \begin{enumerate}
  \def\labelenumii{\alph{enumii})}
  \tightlist
  \item
    Let A be a Noetherian local ring, and let p be a prime ideal of A. Prove the inequality \(\dim_{\mathfrak{p}}(\mathbf{A}) \geqslant \mathrm{ht}(\mathfrak{p}) + \dim(\mathbf{A}/\mathfrak{p}) - 1\) .
  \end{enumerate}
\end{enumerate}

\begin{enumerate}
\def\labelenumi{\alph{enumi})}
\setcounter{enumi}{1}
\tightlist
\item
  Prove that there exists a Noetherian local ring R of dimension 3 having a prime ideal q such that ht(q) = 1 and dim(R/q) = 1. (Use exerc. 16, p.~83, by taking \(\mathbf{R} = \mathbf{A}[\mathbf{U}]_{\mathbf{q}}\) , and q = pR.) Let \(f \in \mathfrak{m}_{\mathbb{R}} - \mathfrak{q}\) ; show that \(\dim(\mathbf{R}_{f}) = 2\) (use exerc. 1, p.~86); we therefore have \(\dim_{\mathfrak{q}}(\mathbf{R}) = 2\) whereas \(\operatorname{ht}(\mathfrak{q}) + \dim(\mathbf{R}/\mathfrak{q}) - 1 = 1.\)
\end{enumerate}

\begin{enumerate}
\def\labelenumi{\arabic{enumi})}
\setcounter{enumi}{3}
\item
  Let \(k\) be a field, and let \(n\) be an integer \(\geqslant 1\). For each integer \(r \in [0, n]\) , let \(M_r\) be the vector subspace of \(k(X_1, ..., X_n)\) generated by the monomials \(X_1^{\alpha_1}\ldots X_r^{\alpha_r}\) with \(\alpha_1, ..., \alpha_r\) in \(Z\) and \(\alpha_r > 0\) . Put \(a_r = M_r + \cdots + M_n\) , for \(0 \leqslant r \leqslant n + 1\) . Then \(a_0\) is a subring of \(k(X_1, ..., X_n)\) , and the \(a_r\) are prime ideals of \(a_0\) . Put \(A = (a_0)_{a_1}\) and \(p_r = (a_r)_{a_1}\) for \(1 \leqslant r \leqslant n + 1\) . Show that the only prime ideals of the local ring \(A\) are the \(p_r\) , and that one has \(ht(p_r) = n + 1 - r\) and \(dim(A) = n\) . Moreover \(p_1\) is principal.
\item
  Let \(\rho: \mathbf{A} \to \mathbf{B}\) be a local homomorphism of complete Noetherian local rings, such that \([\kappa(\mathbf{B}): \kappa(\mathbf{A})] < \infty\) . Put \(n = \dim(\mathbf{B}/\mathfrak{m}_{\mathbf{A}}\mathbf{B})\) ; show that there exists a local A-homomorphism \(u: \mathbf{A}[[\mathbf{T}_1, ..., \mathbf{T}_n]] \to \mathbf{B}\) which makes B a \(\mathbf{A}[[\mathbf{T}_1, ..., \mathbf{T}_n]]\) -finite algebra. (Lift in B a maximal secant sequence of \(\mathbf{B}/\mathfrak{m}_{\mathbf{A}}\mathbf{B}\) and use III, § 3, exerc. 18.)
\item
  Let \(k\) be a field and A the ring of formal power series \(k[[X_1, X_2]]\) . Consider an ideal \(\mathfrak{n}\) of A containing a power of the maximal ideal \(\mathfrak{m}\) of A, and the ideal I generated by \((X_1 - X_2)\mathfrak{n}\) . Let M be the module A/I. Show that the sequence \((X_1)\) is secant for M without being completely secant for M.
\item
  Let A be a local, integral, Noetherian, non-catenary ring of dimension 3 (exerc. 16, p.~83), p ⊂ A a prime ideal such that ht(p) = 1, dim(A/p) = 1. Show that there exists an element x ∈ p such that p is minimal among the prime ideals containing Ax. Deduce that V(x) has an irreducible component of dimension 1.
\end{enumerate}

